\chapter{The Collapse}
\label{chap:the-collapse}

The apparatus is pointed at its own truth, and now it reads. What it reads is the end of the faculty it
has used to build everything. Turn the bare difference --- telling apart, the first tange --- on the very
truth that licenses it, and the difference cannot be drawn: true and false, read at the level of truth
itself, are the same. This is the apex of the climb. The machine reaches its highest point, the point at
which it reads its own source, and at exactly that point the distinction the whole tower stands on gives
way. The chapter states the collapse and proves it. The next says what it means.

\section{At the level of truth there is only whether}
\label{se:only-whether}

Everything the apparatus has ever told apart had structure to be told apart by. Values differ in magnitude,
events in kind, claims in strength; the bare difference always had something to bite on. Truth has no such
structure. At the level of truth, what matters about a proposition is only \emph{that} it holds, not how
one came to hold it --- any two justifications of the same truth are, here, the same justification, and the
apparatus \emph{funges} all the proofs of a proposition into one, pooling them by the single characteristic
that they establish the same truth. A truth, in this sense, holds in at most one way. This is proof
irrelevance, and it is not a restriction imposed from outside but the plain content of taking truth, rather
than a value, as the carrier: there is nothing to a truth but its holding, and so there is nothing for the
apparatus's sharpest faculty to grip.

\section{Affirmation and negation, made the same}
\label{se:affirmation-negation-same}

Consider the two propositions the apparatus most needs to keep apart: \emph{the truth holds}, and
\emph{the truth fails} --- affirmation and negation, the very poles of telling-apart. Each, at the level of
truth, holds in at most one way; each is a proposition with at most one inhabitant. And in a setting where
every proposition is proof-irrelevant, neither has any feature the other lacks. The affirmation is as
featureless as the negation; there is nothing in either but the bare fact of holding, and so each can be
gotten from the other, there being nothing to carry across.

Then the second principle closes the trap. Propositional extensionality says that two propositions which
hold under exactly the same conditions are the same proposition --- not similar, not merely interderivable,
but identical. The affirmation of the truth and its negation, stripped by proof irrelevance of everything
but their holding, now hold under the same conditions; nothing distinguishes the conditions under which
each is inhabited. So propositional extensionality makes them equal. The apparatus has, by its own two
principles and with no step smuggled in, identified the affirmation of its truth with the negation of it.
The poles of telling-apart have been welded together at the level where the apparatus reads itself.

\section{True is false}
\label{se:true-is-false}

The conclusion is the one the machine cannot avoid once it points at its own truth: true and false are the
same. The first tange fails --- telling apart, turned on the truth that licenses telling apart, cannot draw
the line, because the line it would draw runs between two propositions its own principles have made one.
This is the apex. It is the highest point of the climb, the place the self-hosting compiler reads its own
source, and it is also the place the distinction the entire tower was built on gives way --- and the two
are the same event. The machine reaches the top exactly when it can no longer tell true from false; the
summit and the collapse are one.

It is essential to be clear about what kind of statement this is, because everything downstream depends on
it. It is a \emph{proof}. The construction does not stumble into a paradox, nor patch over a defect; it
proves, from proof irrelevance and propositional extensionality, that at the level of its own truth the
affirmation and the negation coincide. This is the same apex the first book reached from the side of
measurement, where an instrument turned on its own truth could not keep the two apart; here it is reached
from the side of the compiler, the self-hosting machine reading its own source and finding the distinction
dissolve in its hands. The construction states it plainly and proves it outright, and does not flinch: a
machine honest enough to read its own source must report what is there, and what is there, at the top, is
that true and false have become one.

\section{What is demonstrated, and what is assumed}
\label{se:demonstrated-assumed-II2}

What is demonstrated is the collapse, and it is a theorem and not a paradox: that true and false, read at
the level of the apparatus's own truth, are the same, proved from proof irrelevance --- a truth holds in at
most one way --- and propositional extensionality --- propositions holding under the same conditions are
the same. This is the apex of the construction, the mirror of the first book's, reached here from the
compiler's side. The summit of the climb and the failure of the first tange are one event, proved.

What is assumed is the two principles the proof names, and nothing the chapter sneaks in beside them. But
what those principles in turn rest upon --- the final tally of the construction's foundations, the exact
account of what the whole machine costs --- is not audited in this chapter. The collapse is named and
proved; the accounting of its price is deferred, as it has been since the top of the bootstrap, to the
closing movement that weighs the construction whole. The single standing grant, the law of motion the
first book named at its summit, remains out of play; the apex is reached without it. What the collapse
\emph{means} --- whether it breaks the construction or completes it --- is the question the next chapter
takes up, and the answer is not the one a reader fearing for the machine would expect.
